% ============================================================================
% Fichier  : partie1/P01_C03_crypto_symetrique.tex
% Titre    : Cryptographie (1) : Symétrique, Hachage, Intégrité
% Auteur   : MINKA MI NGUIDJOÏ Thierry Emmanuel
% Version  : 1.0 -- novembre 2026
% ============================================================================
% Sources : Kerckhoffs (1883) ; FIPS 197 (AES) ; NIST SP 800-38A et 38D
% (modes) ; FIPS 180-4 et FIPS 202 (SHA-2, SHA-3) ; RFC 2104 (HMAC) ;
% NIST SP 800-63B (mémorisation des secrets) ; Wang et Yu (2005), Stevens et
% al. (2017) pour les collisions MD5 et SHA-1 ; loi n° 2010/012, art. 4,
% point 31. Empreintes de l'encadré Outils calculées le 07/11/2026.
% ============================================================================

\chapter{Cryptographie (1) : Symétrique, Hachage, Intégrité}
\label{chap:crypto-sym}
\chaptermark{Cryptographie symétrique}

\epigraph{%
    « Il faut que le système n'exige pas le secret,\\
    et qu'il puisse sans inconvénient tomber entre les mains de l'ennemi. »%
}{Auguste Kerckhoffs, 1883}

\niveauChapitre{Fondamental}{Chapitre~\ref{chap:proprietes}~; notions
d'arithmétique binaire.}

\begin{boxobjectifs}
\begin{itemize}
    \item Énoncer le principe de Kerckhoffs et le vocabulaire de la
          cryptographie.
    \item Expliquer le chiffrement symétrique, par flot et par blocs, et le
          rôle des modes d'opération.
    \item Décrire les propriétés d'une fonction de hachage et ses usages.
    \item Distinguer empreinte, code d'authentification de message (MAC) et
          signature.
    \item Stocker correctement des mots de passe.
    \item Dire ce que la cryptographie ne garantit pas.
\end{itemize}
\end{boxobjectifs}

\bigskip

% ============================================================================
\section*{Question directrice : le chiffrement aurait-il protégé les données volées~?}
% ============================================================================

À Garoua, l'attaquant a copié 40~Go de dossiers clients. S'ils avaient été
chiffrés, la copie aurait-elle été inutile~? La réponse dépend de ce qui
était chiffré, avec quoi, et surtout de l'endroit où se trouvait la clé. Ce
chapitre donne les outils pour répondre.

\medskip
\textbf{Plan du chapitre.}
\begin{enumerate}
    \item Le vocabulaire et le principe de Kerckhoffs (section~\ref{sec:c03-vocab}).
    \item Le chiffrement symétrique (section~\ref{sec:c03-sym}).
    \item Les modes d'opération (section~\ref{sec:c03-modes}).
    \item Les fonctions de hachage (section~\ref{sec:c03-hash}).
    \item Codes d'authentification et mots de passe (section~\ref{sec:c03-mac}).
    \item Ce que la cryptographie ne garantit pas (section~\ref{sec:c03-limites}).
\end{enumerate}

\separateur

% ============================================================================
\section{Le Vocabulaire et le Principe de Kerckhoffs}
\label{sec:c03-vocab}
% ============================================================================

\begin{boxjuris}
\textbf{\loicam{n°~2010/012}, article~4, point~31} (texte officiel)~:
\textit{« Cryptographie~: application des mathématiques permettant d'écrire
l'information, de manière à la rendre inintelligible à ceux ne possédant
pas les capacités de la déchiffrer. »}
\end{boxjuris}

\begin{boxdef}
\termecle{Chiffrement}~: transformation d'un \emph{texte clair} en
\emph{texte chiffré} au moyen d'un algorithme et d'une \emph{clé}.
\termecle{Déchiffrement}~: opération inverse, avec la clé. Le
\emph{décryptement} désigne la retrouvaille du clair \emph{sans} la clé~:
c'est l'attaque.
\end{boxdef}

En 1883, A.~Kerckhoffs pose qu'un système de chiffrement doit rester sûr
même si l'adversaire connaît l'algorithme~\cite{kerckhoffs1883}~: seule
la clé est secrète. C'est le principe de conception ouverte du
Chapitre~\ref{chap:proprietes}. Un algorithme « maison », gardé secret,
n'apporte aucune garantie~; un algorithme public, étudié depuis des années
par la communauté, en apporte.

% ============================================================================
\section{Le Chiffrement Symétrique}
\label{sec:c03-sym}
% ============================================================================

Dans le chiffrement \termecle{symétrique}, la même clé sert à chiffrer et à
déchiffrer. Il est rapide et convient aux grands volumes~; sa difficulté
est de partager la clé (Chapitre~\ref{chap:crypto-asym}).

\begin{tableaucours}{L{3.0cm}L{5.0cm}L{4.2cm}}{%
    \tableauHeader{Algorithme} &
    \tableauHeader{Caractéristiques} &
    \tableauHeader{Statut}
}
DES (1977) & Blocs de 64~bits, clé de 56~bits & Abandonné~: clé trop
courte \\
3DES & Trois passes de DES & Retiré par le NIST \\
AES (2001) & Blocs de 128~bits, clés de 128, 192 ou 256~bits & Norme
actuelle (FIPS~197)~\cite{fips197} \\
ChaCha20 & Chiffrement par flot, rapide en logiciel & Recommandé, utilisé
dans TLS~1.3 \\
\end{tableaucours}
\vspace{-0.8em}
\begin{center}{\small\itshape Tableau~3.1. Algorithmes symétriques}\end{center}

\subsection{La force d'une clé}

Une clé de \(n\)~bits offre \(2^n\) possibilités. À un billion
(\(10^{12}\)) d'essais par seconde, l'espace de \(2^{56}\) clés du DES se
parcourt en une vingtaine d'heures~; celui de \(2^{128}\) clés d'AES
demanderait environ \(10^{19}\)~années. La force brute n'est pas le
problème~: les attaques réelles visent la clé là où elle est stockée, ou
le mauvais usage de l'algorithme.

% ============================================================================
\section{Les Modes d'Opération}
\label{sec:c03-modes}
% ============================================================================

AES chiffre des blocs de 16~octets. Pour chiffrer un fichier, il faut un
\termecle{mode d'opération}, qui dit comment enchaîner les
blocs~\cite{nist80038a}.

\begin{tableaucours}{L{2.4cm}L{6.0cm}L{3.8cm}}{%
    \tableauHeader{Mode} &
    \tableauHeader{Principe} &
    \tableauHeader{À retenir}
}
ECB & Chaque bloc chiffré indépendamment & À proscrire~: deux blocs
identiques donnent le même chiffré \\
CBC & Chaque bloc combiné avec le chiffré précédent, avec un vecteur
d'initialisation & Confidentialité seule, vecteur imprévisible \\
CTR & Chiffrement d'un compteur, combiné au clair & Parallélisable~;
jamais deux fois le même compteur \\
GCM & CTR plus un code d'authentification & Chiffrement
\emph{authentifié}~: confidentialité et intégrité~\cite{nist80038d} \\
\end{tableaucours}
\vspace{-0.8em}
\begin{center}{\small\itshape Tableau~3.2. Les modes d'opération}\end{center}

\begin{boxalert}
\textbf{Violation typique.} Chiffrer sans authentifier (ECB ou CBC seuls).

\smallskip
\textbf{Règle~:} un chiffré non authentifié peut être modifié sans être
lu~; un attaquant change des bits du montant d'un virement sans connaître
la clé. On utilise un chiffrement authentifié (AES-GCM, ChaCha20-Poly1305).
\end{boxalert}

% ============================================================================
\section{Les Fonctions de Hachage}
\label{sec:c03-hash}
% ============================================================================

\begin{boxdef}
\termecle{Fonction de hachage cryptographique}~: fonction qui associe à
un message de taille quelconque une \termecle{empreinte} de taille fixe,
avec trois propriétés~: résistance à la \emph{préimage} (retrouver un
message à partir de l'empreinte est impraticable), à la \emph{seconde
préimage} (trouver un autre message de même empreinte est impraticable),
et aux \emph{collisions} (trouver deux messages de même empreinte est
impraticable).
\end{boxdef}

\begin{tableaucours}{L{2.8cm}L{2.2cm}L{7.2cm}}{%
    \tableauHeader{Fonction} &
    \tableauHeader{Empreinte} &
    \tableauHeader{Statut}
}
MD5 & 128~bits & Collisions trouvées en 2004-2005~\cite{wang2005}~;
interdit pour la sécurité \\
SHA-1 & 160~bits & Première collision publique en 2017~\cite{stevens2017}~;
à abandonner \\
SHA-256, SHA-512 & 256, 512~bits & Famille SHA-2 (FIPS~180-4)~\cite{fips1804}~;
recommandée \\
SHA-3 & 224 à 512~bits & FIPS~202~\cite{fips202}~; alternative de
conception différente \\
\end{tableaucours}
\vspace{-0.8em}
\begin{center}{\small\itshape Tableau~3.3. Fonctions de hachage}\end{center}

\begin{boxoutils}{Une empreinte change entièrement}
\begin{lstlisting}[language=Python]
import hashlib
for m in [b"Virement 500000 FCFA", b"Virement 5000000 FCFA"]:
    print(hashlib.sha256(m).hexdigest())
# 6c82206c2a2b225f5f3ca81065c406972d4e95f04fbdac861cb49d16ce8b9393
# da51d0470e5547e7e81c3cde35325d5514ecf71a5c94553c9e4728f1131639d1
\end{lstlisting}
Un seul zéro ajouté change toute l'empreinte. En ligne de commande~:
\texttt{sha256sum fichier} calcule l'empreinte d'un fichier, et
\texttt{sha256sum -c MANIFESTE} vérifie une liste d'empreintes.
\end{boxoutils}

% ============================================================================
\section{Codes d'Authentification et Mots de Passe}
\label{sec:c03-mac}
% ============================================================================

\subsection{Empreinte, MAC, signature}

Une empreinte seule ne protège pas l'intégrité contre un attaquant~: il
modifie le message et recalcule l'empreinte. Un \termecle{code
d'authentification de message} (MAC) fait intervenir une clé secrète~:
sans elle, impossible de produire un code valide. HMAC en est la
construction standard~\cite{rfc2104}.

\begin{tableaucours}{L{2.8cm}L{2.6cm}L{3.4cm}L{3.4cm}}{%
    \tableauHeader{Mécanisme} &
    \tableauHeader{Clé} &
    \tableauHeader{Protège contre} &
    \tableauHeader{Non-répudiation~?}
}
Empreinte & Aucune & L'erreur accidentelle & Non \\
MAC (HMAC) & Secrète, partagée & La modification par qui n'a pas la clé &
Non~: les deux parties ont la clé \\
Signature & Privée et publique & La modification, et l'usurpation & Oui
(Chapitre~\ref{chap:crypto-asym}) \\
\end{tableaucours}
\vspace{-0.8em}
\begin{center}{\small\itshape Tableau~3.4. Empreinte, MAC et signature}\end{center}

\subsection{Stocker des mots de passe}

Un mot de passe ne se stocke ni en clair, ni chiffré, ni haché par une
simple fonction rapide comme SHA-256~: un attaquant qui vole la base
essaie des milliards de mots de passe par seconde. On utilise une fonction
\emph{lente} et \emph{salée} (Argon2, bcrypt, PBKDF2)~: le \termecle{sel},
aléatoire et propre à chaque compte, empêche les tables précalculées, et la
lenteur réduit le nombre d'essais. Le NIST décrit ces exigences, et
recommande des mots de passe longs plutôt que des changements
fréquents~\cite{nist80063b}.

% ============================================================================
\section{Ce que la Cryptographie ne Garantit pas}
\label{sec:c03-limites}
% ============================================================================

\begin{boxscenario}{Opération GAROUA. Des données chiffrées, mais lisibles}
\textbf{Constatation.} Les dossiers clients étaient stockés sur un
volume chiffré du serveur de fichiers. L'attaquant, devenu administrateur
du domaine, les a copiés \emph{depuis le serveur en fonctionnement}~: le
système les déchiffrait pour lui comme pour tout utilisateur autorisé.

\medskip
\textbf{Analyse} (confiance élevée). Le chiffrement du disque protège
contre le vol du disque, pas contre un attaquant qui a les droits
d'administrateur sur la machine allumée. Pour protéger contre ce
scénario, il fallait limiter les droits (Chapitre~\ref{chap:iam}) et
détecter la copie (Chapitre~\ref{chap:operations}).
\end{boxscenario}

\begin{itemize}
    \item La cryptographie protège des données, pas un système dont
          l'attaquant a le contrôle.
    \item Elle vaut ce que vaut la protection de ses clés.
    \item Elle ne dit rien de l'exactitude de ce qui a été chiffré~: une
          donnée fausse, signée, reste fausse.
\end{itemize}

\begin{boxprocedure}
\textbf{Gérer une clé de chiffrement}
\begin{enumerate}
    \item \textbf{Générer} la clé avec un générateur aléatoire
          cryptographique, jamais à partir d'un mot.
    \item \textbf{Stocker} la clé hors du système qu'elle protège (module
          matériel, coffre de clés).
    \item \textbf{Restreindre} son usage aux seuls services qui en ont
          besoin, et le journaliser.
    \item \textbf{Renouveler} la clé à échéance et en cas de soupçon.
    \item \textbf{Sauvegarder} la clé à part~: une clé perdue rend les
          données irrécupérables.
\end{enumerate}
\end{boxprocedure}

\begin{boxmetier}{Développeur}
\textbf{Ne jamais inventer sa cryptographie.} Le développeur utilise une
bibliothèque reconnue, un chiffrement authentifié, un générateur
aléatoire cryptographique, et une fonction lente et salée pour les mots
de passe. Il ne laisse aucune clé dans le code source.
\end{boxmetier}

\separateur

\begin{boxresume}
\textbf{Ce qu'il faut retenir de ce chapitre~:}
\begin{itemize}
    \item \textbf{Principe de Kerckhoffs}~: seule la clé est secrète.
    \item \textbf{AES} est la norme symétrique~; on l'emploie dans un mode
          \emph{authentifié} (GCM).
    \item \textbf{Une fonction de hachage} résiste à la préimage, à la
          seconde préimage et aux collisions~; MD5 et SHA-1 sont cassés.
    \item \textbf{Empreinte, MAC, signature} n'offrent pas les mêmes
          garanties~; seule la signature apporte la non-répudiation.
    \item \textbf{Les mots de passe} se stockent avec une fonction lente et
          salée.
    \item \textbf{La cryptographie} ne protège pas un système dont
          l'attaquant a pris le contrôle.
\end{itemize}
\end{boxresume}

\begin{boxinfo}
\textbf{Points de contrôle de l'élève.} Je sais~: énoncer le principe de
Kerckhoffs~; expliquer pourquoi ECB est à proscrire~; citer les trois
propriétés d'une fonction de hachage~; distinguer empreinte, MAC et
signature~; expliquer le rôle du sel~; dire ce que le chiffrement de
disque protège et ne protège pas.
\end{boxinfo}

\bigskip

\begin{boxexo}[title={\ding{45}\quad Exercice~3.1. Vocabulaire \niveaubase}]
Distinguez chiffrer, déchiffrer et décrypter, avec un exemple pour chacun.
\end{boxexo}

\begin{boxexo}[title={\ding{45}\quad Exercice~3.2. Kerckhoffs \niveaubase}]
Un éditeur vante un algorithme de chiffrement « secret, donc
inviolable ». Que répondez-vous~?
\end{boxexo}

\begin{boxexo}[title={\ding{45}\quad Exercice~3.3. Force brute \niveaubase}]
Combien de fois plus de clés une clé de 128~bits offre-t-elle qu'une clé
de 56~bits~? Exprimez le rapport en puissance de~2.
\end{boxexo}

\begin{boxexo}[title={\ding{45}\quad Exercice~3.4. ECB \niveaumoyen}]
Une image en noir et blanc est chiffrée avec AES en mode ECB. Pourquoi
peut-on encore deviner son contenu~? Quel mode recommandez-vous~?
\end{boxexo}

\begin{boxexo}[title={\ding{45}\quad Exercice~3.5. Empreinte ou MAC \niveaumoyen}]
LiZbiZ transmet chaque soir le fichier des paiements à sa banque, avec son
empreinte SHA-256 dans le même courriel. Quelle attaque cette protection
n'empêche-t-elle pas~? Que proposez-vous~?
\end{boxexo}

\begin{boxexo}[title={\ding{45}\quad Exercice~3.6. Mots de passe \niveaumoyen}]
La base de l'application de LiZbiZ stocke les mots de passe hachés avec
SHA-256, sans sel. Décrivez l'attaque possible après un vol de la base,
et la correction.
\end{boxexo}

\begin{boxexo}[title={\ding{45}\quad Exercice~3.7. Vérifier un lot \niveaumoyen}]
Écrivez, en Python, une fonction qui reçoit un dossier et un manifeste
d'empreintes SHA-256, et renvoie la liste des fichiers modifiés ou
manquants.
\end{boxexo}

\begin{boxexo}[title={\ding{45}\quad Exercice~3.8. Le chiffrement de disque \niveaumoyen}]
Pour chacun des scénarios, dites si le chiffrement du disque protège les
données~: vol d'un ordinateur portable éteint~; attaquant administrateur
du serveur allumé~; disque jeté sans effacement~; logiciel malveillant
sur le poste de l'utilisateur.
\end{boxexo}

\begin{boxexo}[title={\ding{45}\quad Exercice~3.9. Une politique \niveauexpert}]
Rédigez, pour LiZbiZ, une politique d'usage de la cryptographie d'une page~:
algorithmes autorisés, modes, longueurs de clés, stockage des mots de
passe, gestion des clés.
\end{boxexo}

\begin{boxexo}[title={\ding{45}\quad Exercice~3.10. Les données volées \niveauexpert}]
Reprenez la question directrice. Proposez une architecture dans laquelle
la copie des dossiers clients par un administrateur du domaine aurait été
inutile, et dites ce qu'elle coûte en exploitation.
\end{boxexo}

\bigskip

\begin{boxinfo}
\textbf{Pour aller plus loin.}
\begin{itemize}
    \item FIPS~197 (AES), FIPS~180-4 (SHA-2) et FIPS~202
          (SHA-3)~\cite{fips197,fips1804,fips202}.
    \item NIST SP~800-38D (GCM) et SP~800-63B (secrets
          mémorisés)~\cite{nist80038d,nist80063b}.
    \item J.~Katz et Y.~Lindell, \emph{Introduction to Modern
          Cryptography}~\cite{katz2020}.
\end{itemize}
\end{boxinfo}

\begin{boxinfo}
\textbf{Transition.} La cryptographie symétrique suppose une clé
partagée. Comment la partager avec un inconnu, et comment prouver qui a
signé~? C'est l'objet de la cryptographie asymétrique
(Chapitre~\ref{chap:crypto-asym}).
\end{boxinfo}
